%%%%%%%%%%%%%%%%%%%%%%%%%%%%%%%%%%%%%%%%%%%%%%%%%%%%%%%%%%%%%%%%%%%%%%%%%%%
\chapter*{Conclusion}
\addcontentsline{toc}{chapter}{Conclusion}
%%%%%%%%%%%%%%%%%%%%%%%%%%%%%%%%%%%%%%%%%%%%%%%%%%%%%%%%%%%%%%%%%%%%%%%%%%%

La conclusion générale doit d'abord synthétiser les points clés du rapport en rappelant brièvement les réalisations principales et les objectifs atteints, en soulignant la valeur ajoutée apportée. Elle doit ensuite ouvrir sur les perspectives futures du projet et les axes d'amélioration potentiels.



\begin{tcolorbox}[
    enhanced,
    title={\textbf{Attention}},
    colback=white,
    colframe=red!75!black,
    colbacktitle=red!85!black,
    coltitle=white,
    fonttitle=\large\bfseries,
    sharp corners,
    boxrule=1pt,
    left=10pt,
    right=10pt,
    top=10pt,
    bottom=10pt,
    arc=0pt
]
N'oubliez pas d'inclure à la fin de votre rapport :

\textbf{Une bibliographie et une webographie complètes comprenant} :
\begin{itemize}
    \item Les références des ouvrages et articles consultés
    \item Les liens des sites web consultés avec leurs dates de consultation
    \item Les noms des auteurs lorsqu'ils sont disponibles
\end{itemize}

\textbf{Important} : Évitez de citer des sources non fiables comme Wikipédia. Privilégiez les sources académiques, les sites officiels des éditeurs de logiciels, les documentations techniques vérifiées et les publications professionnelles reconnues.
\end{tcolorbox}